\paragraph{3. Proof of the separator Theorem 14.}
First we need to understand what it means to have slightly non-trivial statistical distance. Let $P$ be a distribution over $\Sigma^m$. One way in which $P$ can have statistical distance $\leq1-\epsilon$ from $S$ is if $P$ is distributed like $S$ with probability $\epsilon$, and it is say fixed with probability $1-\epsilon$. In this case, $P$ has actually very high entropy $(\log_2|S|)$ conditioned on an event of probability $\epsilon$. The next lemma shows that this in fact always happens.
\paragraph{Lemma 15.}
Let $P$ be a distribution over $\Sigma^m$ and let $S\subseteq\Sigma^m$. Suppose that $\Delta(P,S)\leq1-\epsilon$, where $\epsilon\geq\sqrt{8/|S|}$. Then there is a subset $S_0\subseteq S$ of probability $\mathbb P_P[S_0]=\Omega(\epsilon)$ such that the distribution $P$ conditioned on $P\in S_0$ has min-entropy $\geq H_\infty(S)-O(\log1/\epsilon)$.
\begin{proof}
We also write $P$ for the random variable distributed according to $P$. Collect all the elements of $S$ in increasing order of mass according to $P$ until right before collecting cumulative mass $\epsilon/2$. Note we don't collect all of $S$, for else $\mathbb P[P\in S]\leq\epsilon/2$ and $\Delta(P,S)\geq1-\epsilon/2$, contradicting the hypothesis.
Let $\beta$ be the mass of the next element of $S$. Let $S_0$ be the collected elements, $S_1$ the rest of $S$, and $T$ the complement of $S$. By definition, $\mathbb P[P\in S_0]<\epsilon/2$, and so $\mathbb P[P\in S_1\cup T]\geq1-\epsilon/2$. Also for every $x\in S_1$ we have $\mathbb P[P=x]\geq\beta$ and so $|S_1|\leq\beta^{-1}$. Combining these bounds with the assumption we have
\[1-\epsilon\geq\Delta(P,S)\geq\mathbb P[P\in S_1\cup T]-\frac{|S_1|}{|S|}\geq1-\epsilon/2-\frac{\beta^{-1}}{|S|}\]
and so $\beta\leq2/(\epsilon|S|)$.
Because we did not include in $S_0$ an element of mass $\beta$, and we only stop when we reach $\epsilon/2$, the mass of $S_0$ is $\geq\epsilon/2-\beta\geq\epsilon/2-2/(\epsilon|S|)$. If $\epsilon\geq\sqrt{8/|S|}$ this mass is at least $\epsilon/4$.
For any $x\in S_0$ using the above bound on $\beta$ we obtain
\[\mathbb P[P=x\mid P\in S_0]=\frac{\mathbb P[P=x]}{\mathbb P[P\in S_0]}\leq\frac{\beta}{\epsilon/4}\leq\frac8{\epsilon^2|S|},\]
as desired.
\end{proof}
The above lemma allows us to ``forget'' about $S$ and focus on $f$. We need to show that we can restrict the input to a large subset such that many output trees are nearly independent. This is the content of the following theorem. To avoid having to think about infinite sets, in the remainder of the proof we set the input to the sampler to $W^s$ for an integer $s$. This is without loss of generality, since obviously any forest of fixed depth can only access a finite number of input words.
We define the (entropy) loss of a subset $D'\subseteq D$ to be $\log_2(|D|/|D'|)$. So if $D'$ contains half the elements of $D$ the loss is one.
\paragraph{Theorem 16.}
There exists an integer $c\geq1$ such that the following holds:
\textbf{Hypothesis:} Let $f:W^s\to\Sigma^m$ be a depth-$q$ forest with word size $w$. Let $\alpha\leq1/c$. Let $t_0,t_1,\ldots$ be a sequence of integers with $t_i\geq t_{i-1}\cdot cqw/\alpha$ for every $i$. Let $D\subseteq W^s$ be a set with loss $\leq t_0$.
\textbf{Conclusion:} There exists $k$, $1\leq k\leq O(q/\alpha)$, $D_k\subseteq D$, and $t_k$ indices $T\subseteq[m]$ such that:
\begin{enumerate}
\item[(1)] The loss of $D_k\subseteq D$ is $\leq t_{k-1}\cdot(qw/\alpha)^2$;
\item[(2)] For every $i,j\in T$ the random variables $f_i(D_k),f_j(D_k)$ are $O(\alpha)$-independent.
\end{enumerate}
Let us first show how this gives the separator Theorem 14.
\begin{proof}[Proof of Theorem 14 from Theorem 16]
We apply Lemma 15 to $P=f(U)$. Given $S_0\subseteq S$ from the lemma, we let $D\subseteq W^s$ be the preimage of $S_0$ according to $f$. By the lemma, $|D|/|W|^s\geq\Omega(2^{-t_0})$, that is, the loss of $D\subseteq W^s$ is $t_0+O(1)$. Moreover, $H_\infty(f(D))\geq\log|S|-O(t_0)$.
We now apply Theorem 16 to this set $D$ and the sequence $t_0+O(1),t_1,t_2\ldots$. We can adjust the constant $c$ so that this satisfies the hypothesis. The theorem gives $D_k\subseteq D$ with loss $\leq t_{k-1}\cdot(qw/\alpha)^2$.
Observe that the support of $f(D_k)$ is contained in $S$, because the support of $f(D)$ is $S_0\subseteq S$ and $D_k\subseteq D$.
To verify the bound on $H_\infty(f(D_k))$, note that
\[\mathbb P[f(D)=x]\geq\mathbb P[f(D_k)=x]|D_k|/|D|.\]
Taking inverses and then logs we obtain
\[\log(1/\mathbb P[f(D)=x])\leq\log(1/\mathbb P[f(D_k)=x])+\log(|D|/|D_k|).\]
The left-hand side is at least $H_\infty(f(D))\geq\log|S|-O(t_0)$. While $\log(|D|/|D_k|)\leq t_{k-1}\cdot(qw/\alpha)^2$. Hence,
\[\log(1/\mathbb P[f(D_k)=x])\geq\log|S|-O(t_0)-t_{k-1}\cdot(qw/\alpha)^2,\]
for any $x$. The result follows.
\end{proof}
\paragraph{3.1. Proof of Theorem 16.}
The main technical lemma is the following one, which is like Theorem 16 but the requirement of independence is replaced by others easier to work with.
\paragraph{Lemma 17.}
Theorem 16 holds if we replace (2) with:
\begin{enumerate}
\item[(2')] for every $i,j\in T$: the probability over $D_k$ that $f_i(D_k)$ and $f_j(D_k)$ don't make all distinct probes is $\leq\alpha$, and
\item[(2'')] there exists a product distribution $R$ over words (that is, the words are independent) such that for every depth-$2q$ tree $g$, $g(D_k)$ and $g(R)$ are $\alpha$-close.
\end{enumerate}
\paragraph{Lemma 18.}
(2') and (2'') in Lemma 17 imply (2) in Theorem 16.
\begin{proof}
Let $X=D_k$. Think of $(f_i(X),f_j(X))$ as the output of the tree $g$ obtained by appending $f_j$ to the leaves of $f_i$. Note that $g$ makes $2q$ probes, possibly repeated. By (2''), there is a product distribution $R$ such that $g(X)$ and $g(R)$ are $\alpha$-close. Also, the probability that $g$ repeats a probe over $X$ is $\alpha$-close to the probability that it repeats it over $R$. Here we use that this probability can be written as the probability that a tree $d$ of depth $2q$ outputs 1, and that the output distributions of $d$ over $X$ and $R$ are $\alpha$-close. (Tree $d$ is obtained from $g$ by replacing any repeated probe along any path with a leaf 1, and any other leaf with 0.)
By this and (2') the probability that $g$ repeats a probe over $R$ is $\leq2\alpha$. Because $R$ is product, as long as probes are not repeated the output distribution does not change if we answer the first $q$ probes with $R^1$ and the next $q$ probes with $R^2$ where $R^1,R^2$ are iid copies of $R$. This shows that $(f_i(X),f_j(X))$ is $O(\alpha)$-close to $(f_i(R^1),f_j(R^2))$. Using again (2''), we can replace each $R^i$ with $X^i$, where $X^1,X^2$ are iid copies of $X$. This gives that $(f_i(X),f_j(X))$ is $O(\alpha)$-close to $(f_i(X^1),f_j(X^2))$. Adjusting constants concludes the proof.
\end{proof}
\paragraph{3.2. Proof of Lemma 17.}
We shall be concerned with inputs in various subsets $X\subseteq D$. If an input word is constant for every $x\in X$ then it needs not be probed but can be ``hardwired'' in the trees. We shall assume that the trees are always simplified accordingly. We denote by $G(x,X)$ the total number of probes made by all trees on input $x$, where the trees are simplified with respect to $X$.
We use the following fixed-set lemma from [21]. We say that distributions $X$ and $Y$ are $\alpha$-indistinguishable by depth-$q$ trees if for any such tree $t$, the statistical distance between $t(X)$ and $t(Y)$ is $\leq\alpha$.
\paragraph{Lemma 19 ([21], Lemma 3.14).}
Let $B\subseteq W^s$ be a subset with loss $\leq b$, where $W=\{0,1\}^w$. There exists $B_1\subseteq B$ and a product distribution $R$ such that $B_1$ and $R$ are $\alpha$-indistinguishable by depth-$2q$ decision trees. Moreover, the loss of $B_1\subseteq W^s$ is $\leq b\cdot O(wq/\alpha)$.
For completeness we include the proof in Appendix A.
We begin by applying this lemma to $D$ obtaining $D_1\subseteq D$ with loss $t_0\cdot O(wq/\alpha)$. This is the beginning of Iteration 1.
Our goal is to show that at the beginning of Iteration $k$ there exists a subset $D_k\subseteq D$ enjoying the following properties:
\begin{enumerate}
\item[(1)] [in Theorem 16] the loss of $D_k$ is $\leq t_{k-1}\cdot(qw/\alpha)^2$,
\item[(2'')] [in Lemma 17] there exists a product distribution $R$ over words such that for every depth-$2q$ tree $g$, $g(D_k)$ and $g(R)$ are $\alpha$-close.
\item[(3)] $\max_{x\in D_k}G(x,D_k)\leq m(q-\alpha(k-1)/4)$.
\end{enumerate}
Note all these hold at the beginning of Iteration 1.
In an iteration, collect as many trees as possible such that for any two of them, the probability over $D_k$ that they intersect probes is $\leq\alpha$. If you have $t_k$, then (2') in Lemma 17 holds as well, concluding the proof.
Otherwise, you have a collection of $t_k$ trees such that any other tree will intersect a probe with one of those $t$ with probability $\geq\alpha$. Now we are going to use this to proceed to the next iteration, i.e., increase the value of $k$ by 1. Because $G$ is non-negative, Property (3) above implies that there can be at most $O(q/\alpha)$ iterations, as desired.
Write $Y$ for the $\leq t_kq$ words probed by the $t_k$ trees in $D_k$. This is done according to a canonical order, and is a valid definition because the first probe of a tree is fixed, the second is fixed once the answer to the first is, and so on.
\paragraph{Support size.}
Let $D_{k,y}$ be the inputs in $D_k$ with $Y=y$. We have
\[\mathbb E_Y[|D_k|/|D_{k,Y}|]=\sum_y\mathbb P[Y=y]\frac{|D_k|}{|D_{k,y}|}=\sum_y1\leq|W|^{t_kq}.\]
By Markov's inequality $\mathbb P_Y[|D_k|/|D_{k,Y}|\geq M]\leq|W|^{t_kq}/M$. And so with probability $\geq1-|W|^{t_kq}/M$ over $Y$ we have $|D_{k,Y}|\geq|D_k|/M$.
\paragraph{Intersection.}
For a tree $f_i$ let $I_i(x)$ equal 1 if on input $x$ tree $f_i$ intersects probes with at least one of the $t_k$ trees collected, and equal 0 otherwise. Note that for every input $x$ and fixing $y$ we have
\[G(x,D_{k,y})\leq G(x,D_k)-\sum_{i\in[m]}I_i(x).\]
Because $\mathbb P_{x\in D_k}[I_i(x)=1]\geq\alpha$ for every $i$, we have
\[\mathbb E_{x\in D_k}\left[\sum_{i\in[m]}I_i(x)\right]\geq\alpha m.\]
Because the inner sum is $\leq m$, by Markov's inequality we have that with probability $\geq\alpha/2$ over the choice of $Y$
\[\mathbb E_{x\in D_{k,Y}}\left[\sum_{i\in[m]}I_i(x)\right]\geq\alpha m/2,\]
and
\[\mathbb E_{x\in D_{k,Y}}G(x,D_{k,Y})\leq\mathbb E_{x\in D_{k,Y}}G(x,D_k)-\alpha m/2.\]
\paragraph{Combining the arguments.}
Selecting $M=2|W|^{t_kq}/\alpha$ above, and by a union bound, there is a value $\bar y$ so that
\[\mathbb E_{x\in D_{k,\bar y}}[G(x,D_{k,\bar y})]\leq\mathbb E_{x\in D_{k,\bar y}}G(x,D_k)-\alpha m/2;\]
and at the same time $|D_{k,\bar y}|\geq|D_k|\cdot\alpha|W|^{-t_kq}/2$. That is, we increase the loss by
\[\leq\log(1/\alpha)+wt_kq+1.\]
Recall that the loss of $D_k$ is $t_{k-1}\cdot(qw/\alpha)^2$.
Note that $D_{k,\bar y}$ is still uniform over its support, since it is $D_k$ conditioned on a particular choice for $\leq t_kq$ words. Even though the words are chosen adaptively in $D_k$, once we condition on a particular value, their locations are fixed.
\paragraph{Reducing $G$ for every input.}
By Markov's inequality,
\begin{align*}
&\mathbb P_{x\in D_{k,\bar y}}\bigl[G(x,D_{k,\bar y})\geq(\mathbb E_{x\in D_{k,\bar y}}G(x,D_k)-\alpha m/2)(1+\alpha/(4q))\bigr]\\
&\hspace{4em}\leq\frac1{1+\alpha/(4q)}\leq1-\alpha/(8q).
\end{align*}
Hence, for $\geq\alpha/(8q)$ fraction of the inputs $x$ in $D_{k,\bar y}$ we have
\begin{align*}
G(x,D_{k,\bar y})&\leq(\mathbb E_{x\in D_{k,\bar y}}G(x,D_k)-\alpha m/2)(1+\alpha/(4q))\\
&\leq m(q-\alpha(k-1)/4)(1+\alpha/(4q))-\alpha m/2\\
&\leq m(q-\alpha k/4),
\end{align*}
using (3) in the second inequality. Let $D'_{k,\bar y}$ be the set of these inputs. The above gives the desired bound on $\max_{x\in D'_{k,\bar y}}G(x,D'_{k,\bar y})$, and note that the loss of $D'_{k,\bar y}\subseteq D_{k,\bar y}$ is $\leq\log(8q/\alpha)$.
\paragraph{Fixed-set lemma.}
Finally, we apply the fixed-set lemma to $D'_{k,\bar y}$ to obtain $D_{k+1}$; this gives (2''). This application multiplies the loss by $O(wq/\alpha)$, bringing the loss of $D_{k+1}\subseteq D$ to
\[O(wq/\alpha)\cdot O(t_{k-1}\cdot(qw/\alpha)^2+\log(1/\alpha)+wt_kq+1).\]
We need this loss to be at most $t_k\cdot(qw/\alpha)^2$. Dividing by $wq/\alpha$ we need to verify that
\[O(t_{k-1}\cdot(qw/\alpha)^2)+O(\log1/\alpha)+O(wt_kq)+O(1)\leq t_k\cdot qw/\alpha.\]
We claim that each term on the left-hand side is at most one-fourth of the right-hand side. For the first term we use the hypothesis that $t_k\geq t_{k-1}\cdot cqw/\alpha$ for a large enough $c$, and for the third we use that $\alpha\leq1/c$ and pick $c$ large enough. This gives (1).
Because $D_{k+1}\subseteq D'_{k,\bar y}$, the bound on $G$ still holds for $D_{k+1}$, and this gives (3).
